\section{The last auxiliary coagulation: tail bounds, daughter envelopes,
and exact critical representations}
\label{sec:last-event}

Fix a solution in \cref{def:solution} and its number process from
\cref{thm:number-process}. Retain the coagulation count $K_t$ of
\cref{thm:controlled-path}, and define
\begin{equation}
 Q_t=\frac{N(t)X_t}{m},\qquad
 L=\sup\{s>0:s\text{ is an auxiliary coagulation time}\},\qquad
 h(t)=\Prob(L>t),
 \label{eq:last-event-def}
\end{equation}
where the supremum of the empty set is zero. Thus $\E Q_t=1$ and
$\E Q_t^p=\mathcal A_p(t)$. Although $L$ is finite almost surely by
\cref{thm:controlled-path}, it is generally not a stopping time, and
nothing below conditions at $L$. All conditional future laws mean the
Markov construction restarted at a deterministic time $t$ and state
$x$ in the prescribed environment $n_s$; this specifies them even at
states of zero probability under $\eta_t$. By the Markov property of
the Poisson-random-measure construction, the regular conditional law
of $(X_s)_{s\geq t}$ given $X_t$ equals this restart law for
$\eta_t$-almost every $x$, so averaging restart quantities against
$\eta_t$ gives unconditional probabilities.

\subsection{Controlled first-event and last-event bounds}

\begin{theorem}[Controlled first-event and last-event bounds]
\label{thm:last-controlled}
For $0\leq t<T\leq\infty$, set
\begin{equation}
 u_{t,T}=1-\exp\left(-\int_t^T b(s)\dd s\right),\qquad
 u_t=u_{t,\infty}.
 \label{eq:remaining-activity}
\end{equation}
Given $X_t=x$, write $q=N(t)x/m$. Then
\begin{equation}
 \Prob(K_T>K_t\mid X_t=x)
 \leq1-e^{-q u_{t,T}}\leq\min(1,q u_{t,T}).
 \label{eq:first-event-bound}
\end{equation}
At $T=\infty$ the event means at least one coagulation after $t$.
For every $0<p<1$,
\begin{align}
 \Prob(K_T>K_t)
 &\leq u_{t,T}^p\mathcal A_p(t)
 \leq u_{t,T}^p\mathcal A_p(0)
       e^{-a_p\Bc(t)-a_{1-p}F(t)},\label{eq:window-tail}\\
 h(t)&\leq u_t^p\mathcal A_p(0)
       e^{-a_p\Bc(t)-a_{1-p}F(t)}.\label{eq:last-tail}
\end{align}
Together with finite-horizon nonexplosion from
\cref{thm:number-process}, these bounds imply $L<\infty$ and
$K_\infty<\infty$ almost surely, for all allowed controls. No
logarithmic or second size moment is required.
\end{theorem}

\begin{proof}
Start a fragmentation-only process $V_s$ at $V_t=x$, with the
prescribed daughter kernel evaluated at its own current size, driven
by the fragmentation clock and uniform innovations of the auxiliary
construction. It has finitely many events on each finite interval,
$0<V_s\leq x$, and conditional daughter mean $V_{s-}/2$, so
compensation and the integrating factor give
\begin{equation}
 \E_{t,x}V_s=xe^{-[F(s)-F(t)]},\qquad
 \E_{t,x}\frac{N(s)V_s}{m}=q e^{-[\Bc(s)-\Bc(t)]}.
 \label{eq:fragmentation-first-mean}
\end{equation}
The cancellation in the second identity is why only remaining
coagulation activity enters \eqref{eq:first-event-bound}.

Until its first coagulation after $t$, the auxiliary process has
state $V_s$ and coagulation rate $\lambda(s)N(s)V_{s-}$. The
coagulation and fragmentation Poisson measures are independent, so
conditional on the whole fragmentation path the retained coagulation
points before the first coagulation form a Poisson process with
intensity $\lambda(s)N(s)V_s\dd s$. The first future coagulation is
therefore the first crossing of an independent mean-one exponential
threshold by the hazard
\begin{equation}
 \mathcal H_{t,T}=\int_t^T b(s)\frac{N(s)V_s}{m}\dd s,
 \label{eq:first-hazard}
\end{equation}
after which a partner is sampled and the full process resumes; the
full future compensator is not identified with \eqref{eq:first-hazard}.
Tonelli's theorem gives
$\E_{t,x}\mathcal H_{t,T}=q\int_t^Tb(s)e^{-[\Bc(s)-\Bc(t)]}\dd s
=qu_{t,T}$, finite even at $T=\infty$. The event probability is
exactly $\E_{t,x}(1-e^{-\mathcal H_{t,T}})$, and Jensen's inequality
for the concave map $z\mapsto1-e^{-z}$ gives
\eqref{eq:first-event-bound}. Strict concavity makes the first inequality
an equality if and only if $\mathcal H_{t,T}=qu_{t,T}$ almost surely.
Vanishing fragmentation on $(t,T)$ is sufficient, but is not necessary:
fragmentation after all coagulation activity in the window has ended
does not affect the hazard. The second inequality is strict whenever
$0<qu_{t,T}<\infty$.
Averaging over $\eta_t$, then using $\min(1,z)\leq z^p$ and
\cref{thm:affinity}, gives the two tail bounds.

If $\Bc(\infty)=\infty$, the coagulation exponent in
\eqref{eq:last-tail} diverges; if $\Bc(\infty)<\infty$, then
$u_t\to0$ and $\mathcal A_p(t)\leq1$. In both cases $h(t)\to0$, so
almost surely there is a finite time after which no coagulation
occurs, and nonexplosion on that finite horizon leaves finitely many
coagulations before it. Hence $L<\infty$ and $K_\infty<\infty$.
\end{proof}

The conditional bound \eqref{eq:first-event-bound} is thus
instantaneously sharp as a uniform bound over the full controlled
class; its averaged consequence need not have a sharp time exponent.
Each right-hand side in \cref{eq:window-tail,eq:last-tail} may be
multiplied by the scalar factor
\begin{equation}
 c_p=\max_{z>0}\frac{1-e^{-z}}{z^p}<1,\qquad c_{1/2}\approx0.638,
 \label{eq:scalar-tail-factor}
\end{equation}
because the ratio tends to zero at both endpoints and its derivative
changes sign exactly at the unique positive root of $z/(e^z-1)=p$.
This is sharp scalar domination, not a model asymptotic assertion.

\begin{corollary}[Approximation of the entire future path]
\label{cor:future-path-TV}
Start a fragmentation-only process $V^{(t)}$ at the random state
$V^{(t)}_t=X_t$, with daughter kernel $b_{s,V^{(t)}_{s-}}/2$.
For $T\leq\infty$, the total variation distance between the laws of
$(X_s)_{t\leq s\leq T}$ and $(V^{(t)}_s)_{t\leq s\leq T}$ on the Borel
space of c\`adl\`ag $E$-valued paths is at most
\begin{equation}
 u_{t,T}^p\mathcal A_p(0)e^{-a_p\Bc(t)-a_{1-p}F(t)}.
 \label{eq:future-path-TV}
\end{equation}
\end{corollary}

\begin{proof}
Use common future fragmentation times and uniform innovations for
both processes; after their states diverge each quantile map uses its
own state, so each marginal keeps its daughter dynamics. The paths
agree on the whole interval unless $K_T>K_t$, so any two path
probabilities differ by at most $\Prob(K_T>K_t)$; apply
\eqref{eq:window-tail}. The coupling is defined for all future times,
which covers $T=\infty$.
\end{proof}

This reference starts from $\eta_t$, not $\eta_0$. For
parent-dependent daughters, common innovations do not mean common
fractions after divergence, consistent with \cref{rem:adaptive-marks}.

\begin{corollary}[Time moments and rare event-bearing paths]
\label{cor:last-moments-counts}
Let $b>0$ and $\sigma\geq0$ be constant, and put
\begin{equation}
 \chi_p=b a_p+\sigma a_{1-p}>0.
 \label{eq:last-rate}
\end{equation}
Then $h(t)\leq\mathcal A_p(0)e^{-\chi_p t}$, and, for $0<r<\chi_p$,
\begin{equation}
 \E e^{rL}\leq1+\frac{r\mathcal A_p(0)}{\chi_p-r}.
 \label{eq:last-exp-moment}
\end{equation}
For a fixed window length $\ell>0$, let
$D_t=K_{t+\ell}-K_t$ and $\omega_p=(1-e^{-b\ell})^p\mathcal A_p(0)$.
Then $\E D_t=b\ell$, $D_t=0$ eventually almost surely, and
\begin{align}
 \Prob(D_t>0)&\leq\omega_p e^{-\chi_p t},\label{eq:rare-window}\\
 \E[D_t\mid D_t>0]&\geq\frac{b\ell}{\omega_p}e^{\chi_p t},
                     \label{eq:rare-mean}\\
 \E D_t^r&\geq(b\ell)^r\omega_p^{1-r}e^{(r-1)\chi_p t}
                     \quad(r>1).\label{eq:rare-power}
\end{align}
The last inequality allows an infinite left-hand side.
\end{corollary}

\begin{proof}
The time-tail bound is \eqref{eq:last-tail} with $u_t=1$, and Tonelli
applied to $e^{rL}-1=\int_0^Lre^{rs}\dd s$ proves
\eqref{eq:last-exp-moment}. The count compensator has mean
$\int_t^{t+\ell}b\E Q_s\dd s=b\ell$, whereas the finite last time
makes $D_t$ eventually zero on each path. The window bound is
\eqref{eq:window-tail}; dividing the positive mean by
$\Prob(D_t>0)$ gives \eqref{eq:rare-mean}, and H\"older gives
$b\ell\leq(\E D_t^r)^{1/r}\Prob(D_t>0)^{1-1/r}$.
\end{proof}

The window counts are not uniformly integrable despite their
almost-sure disappearance: their mean concentrates on rare auxiliary
paths. No burst-size law, reactor count, or physical mass-tag event
is asserted.

\subsection{A classical pure-coagulation benchmark}
\label{sec:borel-benchmark}

For monodisperse data $n_0=N_0\delta_{x_0}$ and pure coagulation,
write $b=\lambda m>0$ and $\alpha=1-e^{-bt}$. The Golovin solution
\eqref{eq:classical-Borel} \citep{Golovin1963}, the marginal of the
additive coalescent \citep{AldousPitman1998,Bertoin2002}, says after
rescaling that $X_t/x_0$ has law $\operatorname{Borel}(\alpha)$, whose
generating function satisfies
$G_\alpha(z)=z\exp\{\alpha(G_\alpha(z)-1)\}$ by the Poisson branching
decomposition of total progeny. Because $N(t)=N_0(1-\alpha)$, the
normalized state is $Q_t=(1-\alpha)X_t/x_0$ and the no-future-event
probability is exactly $e^{-Q_t}$. Hence
\begin{equation}
 h(t)=1-G_\alpha(e^{-(1-\alpha)}),\qquad
 -\log(1-h(t))=(1-\alpha)+\alpha h(t).
 \label{eq:Borel-last-exact}
\end{equation}
This is a deduction from the Golovin--Borel solution as recorded by
\citet[equation~(3)]{Bertoin2009}, not a new solution. Since
$h(t)\to0$, subtracting $h(t)$ in \eqref{eq:Borel-last-exact} yields
$h(t)^2(1/2+O(h(t)))=(1-\alpha)(1-h(t))$, and therefore
\begin{equation}
 h(t)\sim\sqrt2\,e^{-bt/2}.
 \label{eq:Borel-last-asymptotic}
\end{equation}
Even this solvable endpoint has a faster tail than the certificate
$e^{-a_{1/2}bt}$ from \cref{cor:last-moments-counts} with $\sigma=0$
and $p=1/2$, where $a_{1/2}=\kappa/2$. Sharpness of a moment
differential inequality does not imply sharpness of its derived
last-event exponent.

\subsection{Daughter-law envelopes at constant rates}
\label{sec:daughter-comparison}

From here on $b=\lambda m>0$ and $\sigma\geq0$ are constant; the
competing daughter kernels retain all the time and parent dependence
allowed in \cref{ass:rates}. For a restart at time $t$ and normalized
state $q$, let $g_t(q)$ be the probability of a future auxiliary
coagulation and $w_t(q)=1-g_t(q)$. By the first-event construction
these do not depend on future partner laws, though they may depend on
$t$ through the daughter kernel.

Set $d=\sigma-b$ and $J_s=\int_0^se^{du}\dd u$. For $\sigma>0$, let
$\tau$ be exponential of rate $\sigma$ and $P$ a Poisson process of
rate $2\sigma$, and define
\begin{align}
 Z_{\rm ext}&=bJ_\tau,&
 Z_{\rm eq}&=b\int_0^\infty e^{ds}2^{-P_s}\dd s,
 \label{eq:benchmark-hazards}\\
 w_{\rm ext}(q)&=\E e^{-qZ_{\rm ext}},&
 w_{\rm eq}(q)&=\E e^{-qZ_{\rm eq}}.
 \label{eq:benchmark-transforms}
\end{align}
For $\sigma=0$ define both hazards to equal one, so both transforms
equal $e^{-q}$. The first benchmark is the hazard of a comparison
process growing at rate $d$ until a reset to zero at rate $\sigma$;
the zero-reset rule is not an admissible daughter law and is never
asserted to be one. The second is the hazard of equal splitting.
Tonelli gives $\E Z_{\rm ext}=\E Z_{\rm eq}=1$, since
$e^{ds}\Prob(\tau>s)=e^{-bs}$ and $e^{ds}\E2^{-P_s}=e^{-bs}$. The
second moment of $Z_{\rm ext}$ is finite exactly when $\sigma<2b$;
only first moments are used below.

At criticality, $d=0$, $Z_{\rm ext}=b\tau$ is a mean-one exponential
variable, and $w_{\rm ext}(q)=(1+q)^{-1}$. For equal splitting the
holding times are independent exponentials of rate $2b$, so with
independent mean-one exponentials $E_k$,
\begin{equation}
 Z_{\rm eq}\overset{d}{=}S:=\sum_{k=1}^\infty2^{-k}E_k,
 \qquad \Psi(q):=\E e^{-qS}
       =\prod_{k=1}^\infty(1+q/2^k)^{-1}.
 \label{eq:critical-product}
\end{equation}
This is the Poissonian exponential functional of
\citet[equations~(1.3), (1.6), Theorem~1.1(i)]{BertoinBianeYor2004},
whose Theorem~1.1(i) is stated for parameter values in $(0,1)$; in
their notation $S$ has law $I^{(1/2)}/2$. The product converges
locally uniformly since $\sum_k\log(1+q/2^k)\leq q$, and independence
of the holding times gives
\begin{equation}
 \E S=1,\qquad \Var S=\tfrac13,\qquad
 (2+q)\Psi(q)=2\Psi(q/2),\qquad
 \Psi(q)\leq\frac1{1+q},
 \label{eq:product-identities}
\end{equation}
the last because $\prod_k(1+q/2^k)\geq1+q\sum_k2^{-k}=1+q$. The
Laplace representation makes $\Psi$ strictly decreasing, strictly
convex, and completely monotone, with $\Psi(q)\to0$ as $q\to\infty$.
A large-$q$ form and a rational certificate for $\Psi(1)$ are in
\cref{app:product}.

\begin{theorem}[Conditional envelopes at all constant rates]
\label{thm:daughter-envelopes}
For every allowed daughter kernel, every deterministic $t\geq0$,
and every $q>0$,
\begin{equation}
 w_{\rm eq}(q)\leq w_t(q)\leq w_{\rm ext}(q),\qquad
 1-w_{\rm ext}(q)\leq g_t(q)\leq1-w_{\rm eq}(q).
 \label{eq:daughter-envelopes}
\end{equation}
At criticality $\sigma=b$ these read
\begin{equation}
 \Psi(q)\leq w_t(q)\leq\frac1{1+q},\qquad
 \frac q{1+q}\leq g_t(q)\leq1-\Psi(q),
 \label{eq:critical-envelopes}
\end{equation}
and for equal splitting $g_t(q)=g(q):=1-\Psi(q)$. Both envelopes are
instantaneously sharp as uniform bounds over the stated class: equal
splitting attains the lower bound on $w_t$, and the positive binary
daughter measures
\begin{equation}
 b^{(\varepsilon)}_{s,x}
       =\delta_{\varepsilon x}+\delta_{(1-\varepsilon)x},
       \qquad 0<\varepsilon<1/2,
 \label{eq:epsilon-daughters}
\end{equation}
approach the upper bound as $\varepsilon\downarrow0$. At $\sigma=0$
the envelopes coincide for every daughter kernel.
\end{theorem}

\begin{proof}
For the fragmentation-only future write
$Y_s=N(t+s)V_{t+s}/m$, $Y_0=q$. Between events it has drift $dY_s$;
at rate $2\sigma$ it is multiplied by a fraction $\Theta$ with
conditional mean $1/2$. By \eqref{eq:fragmentation-first-mean},
\begin{equation}
 \E Y_s=qe^{-bs},\qquad
 \mathcal H_s=\int_0^s bY_v\dd v,\qquad
 w_t(q)=\E e^{-\mathcal H_\infty},\qquad \E\mathcal H_\infty=q.
 \label{eq:constant-first-hazard}
\end{equation}
Both benchmark hazards are strictly positive with mean one, so their
transforms $w$ are decreasing, convex, and continuously differentiable
on $[0,\infty)$ (from the right at zero, by dominated convergence
using the finite mean), with
\begin{equation}
 w(0)=1,\qquad |w'(q)|\leq1,\qquad 0\leq1-w(q)\leq q.
 \label{eq:benchmark-regularity}
\end{equation}

\emph{Benchmark equations.} For $\sigma>0$ and $q\geq0$,
\begin{align}
 dq w_{\rm ext}'(q)+\sigma[1-w_{\rm ext}(q)]
                  -bq w_{\rm ext}(q)&=0,\label{eq:ext-benchmark-ODE}\\
 dq w_{\rm eq}'(q)+2\sigma[w_{\rm eq}(q/2)-w_{\rm eq}(q)]
                  -bq w_{\rm eq}(q)&=0;\label{eq:eq-benchmark-ODE}
\end{align}
for $\sigma=0$ both hold with $d=-b$ and $w=e^{-q}$. For the first,
$w_{\rm ext}(q)=\int_0^\infty\sigma e^{-\sigma s}e^{-qbJ_s}\dd s$, and
$\E Z_{\rm ext}=1$ justifies differentiation under the integral.
Since $\partial_se^{-qbJ_s}=-qb(1+dJ_s)e^{-qbJ_s}$ and
$e^{-\sigma s}e^{-qbJ_s}\to0$ as $s\to\infty$, integration by parts
gives
\[
 1-w_{\rm ext}(q)
 =qb\int_0^\infty e^{-\sigma s}(1+dJ_s)e^{-qbJ_s}\dd s
 =\frac{qb}\sigma w_{\rm ext}(q)-\frac{qd}\sigma w_{\rm ext}'(q),
\]
which is \eqref{eq:ext-benchmark-ODE}. For the second, let $\tau_1$
be the first jump time of $P$, exponential of rate $2\sigma$. The
strong Markov property of $P$ at $\tau_1$ gives
$Z_{\rm eq}=bJ_{\tau_1}+e^{d\tau_1}Z'_{\rm eq}/2$ in law, with
$Z'_{\rm eq}$ a copy of $Z_{\rm eq}$ independent of $\tau_1$. Hence
\[
 w_{\rm eq}(q)=\int_0^\infty2\sigma e^{-2\sigma s}e^{-qbJ_s}
                w_{\rm eq}(qe^{ds}/2)\dd s.
\]
If $d=0$ this reads $(2\sigma+bq)w_{\rm eq}(q)=2\sigma w_{\rm eq}(q/2)$,
which is \eqref{eq:eq-benchmark-ODE} and the recurrence in
\eqref{eq:product-identities}. If $d\neq0$, substitute $r=qe^{ds}$:
\[
 w_{\rm eq}(q)=\frac{2\sigma}{|d|}\,q^{2\sigma/d}e^{bq/d}
   \int_{I_q}r^{-2\sigma/d-1}e^{-br/d}w_{\rm eq}(r/2)\dd r,
\]
with $I_q=(q,\infty)$ for $d>0$ and $I_q=(0,q)$ for $d<0$.
Differentiating this product by the fundamental theorem of calculus,
$w_{\rm eq}$ being continuous, gives
$dq\,w_{\rm eq}'(q)=(2\sigma+bq)w_{\rm eq}(q)-2\sigma w_{\rm eq}(q/2)$,
which is \eqref{eq:eq-benchmark-ODE}.

\emph{Comparison.} For any convex $w$ with $w(0)=1$ and every
admissible conditional fraction law, Jensen's inequality and the
chord on $[0,q]$ give
\begin{equation}
 w(q/2)\leq\E[w(q\Theta)\mid s,q]
               \leq\tfrac12[1+w(q)].
 \label{eq:convex-daughter-chord}
\end{equation}
Let $\mathcal G^{\rm frag}_sf(q)=dqf'(q)
+2\sigma\{\E[f(q\Theta)\mid t+s,q]-f(q)\}$ be the fragmentation-only
generator in the $Y$ coordinate, with the actual conditional fraction
law at absolute time $t+s$. The benchmark equations and
\eqref{eq:convex-daughter-chord} give
\[
 (\mathcal G^{\rm frag}_s-bq)w_{\rm ext}(q)\leq0,\qquad
 (\mathcal G^{\rm frag}_s-bq)w_{\rm eq}(q)\geq0.
\]
Dynkin's formula applies to $e^{-\mathcal H_s}w(Y_s)$ for either
benchmark transform: the process lies in $[0,1]$; between jumps
$Y_s>0$ follows the deterministic flow and $w\in C^1$ on $[0,\infty)$
with $|w'|\leq1$ and $Y_s\leq qe^{d_+s}$, $d_+=\max(d,0)$, so the
drift term is integrable on finite horizons; and
$\E\int_0^s2\sigma|w(Y_v\Theta)-w(Y_v)|\dd v\leq2\sigma s<\infty$
because $0\leq w\leq1$. Hence $e^{-\mathcal H_s}w_{\rm ext}(Y_s)$ is a
supermartingale and $e^{-\mathcal H_s}w_{\rm eq}(Y_s)$ a submartingale.
By \eqref{eq:benchmark-regularity},
$\E|e^{-\mathcal H_s}w(Y_s)-e^{-\mathcal H_s}|\leq\E Y_s=qe^{-bs}\to0$,
while $e^{-\mathcal H_s}\to e^{-\mathcal H_\infty}$ boundedly. Taking
expectations and letting $s\to\infty$ proves
\eqref{eq:daughter-envelopes}; no monotonicity of normalized size,
stationarity of the competing kernel, or daughter log moment is used.
The critical forms \eqref{eq:critical-envelopes} are the case $d=0$
with the benchmarks computed above.

\emph{Sharpness.} For equal splitting, $Y_s=qe^{ds}2^{-P_s}$, so the
lower bound on $w_t$ is attained. For \eqref{eq:epsilon-daughters},
use independent clocks of rate $\sigma$ for the two marks, common for
all $\varepsilon$. Let $\tau_\varepsilon$ be the first
$\varepsilon$-mark and $P'_s$ the count of the other marks. The hazard
before $\tau_\varepsilon$ is
\[
 \mathcal H^{\rm pre}_\varepsilon
 =bq\int_0^{\tau_\varepsilon}e^{ds}(1-\varepsilon)^{P'_s}\dd s
       \longrightarrow qZ_{\rm ext}\quad\hbox{in }L^1,
\]
by pathwise convergence and domination by $qZ_{\rm ext}$. Immediately
after $\tau_\varepsilon$,
$Y_{\tau_\varepsilon}=q\varepsilon e^{d\tau_\varepsilon}
(1-\varepsilon)^{P'_{\tau_\varepsilon}}$, and the conditional expected
remaining hazard is $Y_{\tau_\varepsilon}$ by
\eqref{eq:constant-first-hazard} and the fresh clocks at this stopping
time. Independence of the clocks gives
\begin{equation}
 \E Y_{\tau_\varepsilon}=q\varepsilon\sigma\int_0^\infty
                   e^{-(b+\sigma\varepsilon)s}\dd s
          =\frac{q\varepsilon\sigma}{b+\sigma\varepsilon}\to0.
 \label{eq:epsilon-residual}
\end{equation}
Thus the full hazard converges in $L^1$ to $qZ_{\rm ext}$, and since
$e^{-z}$ is one-Lipschitz on $[0,\infty)$ the survival probability
tends to $w_{\rm ext}(q)$. When $\sigma=0$ the hazard is
deterministically $q$.
\end{proof}

Convex inequalities with generator verification are established
comparison methods; ordered-generator comparisons for
time-inhomogeneous Markov processes appear in
\citet{RuschendorfSchnurrWolf2016}. The conclusion orders future-event
probabilities; it asserts neither convex order of hazards nor an
order of total coagulation counts.

\begin{remark}[Exponential functionals of subordinators]
\label{rem:exponential-functional}
At criticality with parent-independent fractions $\Theta$,
$Y_s=qe^{-\xi_s}$, where $\xi$ is the compound-Poisson subordinator
with jump rate $2b$ and jump law that of $-\log\Theta$. The future
hazard $qb\int_0^\infty e^{-\xi_s}\dd s$ is an exponential functional
of a L\'evy process in the sense of \citet{BertoinYor2005}, and $w_t$
is its Laplace transform. This is why the Poissonian exponential
functional of \citet{BertoinBianeYor2004} appears for equal
splitting, where $\xi$ is $\log2$ times a Poisson process of rate
$2b$ and $b\int e^{-\xi_s}\dd s$ has the law of $S$. In this language
\cref{thm:daughter-envelopes} compares $\E e^{-qbI_\xi}$ over jump
laws of $\xi$ with $\E e^{-\Delta\xi}=1/2$: it is the convex-order
statement \eqref{eq:convex-daughter-chord} between the deterministic
fraction $1/2$, the law of $\Theta$, and the degenerate limit
concentrated on $\{0,1\}$.
\end{remark}

\begin{corollary}[Overlap coefficients at all constant rates]
\label{cor:overlap-all-rates}
Write
\begin{equation}
 O(t)=1-\norm{\eta_t-\pi_t}_{\TV}
     =\int_E\min\left(1,\frac{N(t)x}m\right)\eta_t(\dd x)
     =\E\min(1,Q_t),\qquad
 c_{b,\sigma}=1-w_{\rm ext}(1)>0.
 \label{eq:overlap-last-def}
\end{equation}
Then
\begin{equation}
 c_{b,\sigma}O(t)\leq h(t)\leq O(t),
 \label{eq:overlap-all-rates}
\end{equation}
with $c_{b,0}=1-e^{-1}$ and $c_{b,b}=1/2$. For equal splitting at
criticality, moreover,
\begin{equation}
 c_*O(t)\leq h(t)\leq O(t),\qquad
 c_*=1-\Psi(1)=0.580577558204892\ldots
 \label{eq:equal-overlap}
\end{equation}
Each lower coefficient and the common upper coefficient one are
instantaneously sharp as uniform constants over the stated class of
daughter kernels and initial populations.
\end{corollary}

\begin{proof}
The overlap identity follows from $\dd\pi_t/\dd\eta_t=Q_t$ and the
convention \eqref{eq:variation}. The increasing concave function
$f(q)=1-w_{\rm ext}(q)$ vanishes at zero, so its chord on $[0,1]$ and
monotonicity on $[1,\infty)$ give $f(q)\geq f(1)\min(1,q)$. Average
the lower envelope in \eqref{eq:daughter-envelopes} over $\eta_t$ for
the lower bound and use \eqref{eq:first-event-bound} with $u_t=1$ for
the upper bound. At criticality $f(1)=1/2$; for equal splitting the
same chord argument applies to $g$, with $g(1)=c_*$.

For the lower constants, monodisperse data give $Q_0=1$ and $O(0)=1$;
equal splitting attains $c_*$, and the $\varepsilon$-daughters
approach $c_{b,\sigma}$. For the upper constant use the mean-one
two-point laws
\begin{equation}
 \nu_{p,s}=(1-p)\delta_s+p\,\delta_{r(p,s)},\qquad
 r(p,s)=\frac{1-(1-p)s}p,\qquad 0<p<1,\ 0<s<1,
 \label{eq:two-point-family}
\end{equation}
where $\delta_x$ is a Dirac measure. Under equal splitting at
criticality with $Q_0\sim\nu_{\varepsilon,\varepsilon^2}$,
$O(0)=(1-\varepsilon)\varepsilon^2+\varepsilon\sim\varepsilon$ and
$h(0)=(1-\varepsilon)g(\varepsilon^2)+\varepsilon g(r)\sim\varepsilon$,
because $g(q)\leq q$ and $g(r)\to1$ as $r\sim\varepsilon^{-1}\to\infty$.
\end{proof}

All sharpness preparations are realizable populations: for any
finitely supported positive mean-one law $\mu$, the measure
$n_0=N_0(q\mapsto mq/N_0)_\#\mu$ has count $N_0$, mass $m$, and finite
second moment, so \cref{thm:existence} supplies a solution with
$Q_0\sim\mu$. The two-sided bounds make vanishing of the overlap
equivalent to disappearance of future coagulation in this auxiliary
representation. These constants compare probabilities at arbitrary
times and populations; they do not identify a decay exponent.

\subsection{Exact critical representations of the last-event law}
\label{sec:critical-last-event}

Assume from here on that $b=\sigma>0$, so $N(t)=N_0$. Let $\mu_t$ be
the law of $Q_t=N_0X_t/m$, and let $Q'_t$ denote an independent
variable with law $\mu_t$ whenever it appears with $Q_t$. The
generator of $Q$ is
\begin{equation}
 \mathcal G_t f(q)=2b\{\E[f(q\Theta)\mid t,q]-f(q)\}
       +bq\int_E[f(q+v)-f(q)]\mu_t(\dd v),
 \label{eq:critical-Q-generator}
\end{equation}
with the original daughter kernel expressed in normalized size in the
first term, so $\E\Theta=1/2$ conditionally. The formulas below are
exact identities in terms of the unknown law $\mu_t$: representations,
not closed-form evaluations.

\begin{theorem}[Equal-split tail, density, and Laplace representation]
\label{thm:exact-last-density}
For equal splitting and $\Psi$ from \eqref{eq:critical-product},
\begin{align}
 h(t)&=1-\int_E\Psi(q)\mu_t(\dd q),\label{eq:exact-last-observable}\\
 h'(t)&=-b\E[Q_t\Psi(Q_t+Q'_t)].\label{eq:exact-last-density}
\end{align}
Here $h'$ is continuous on $[0,\infty)$, the value at zero being a
right derivative. The law of $L$ has an atom $\E\Psi(Q_0)$ at zero and
density $b\E[Q_t\Psi(Q_t+Q'_t)]$ on $(0,\infty)$. For the size Laplace
transform $\phi_t(z)=\int_Ee^{-zq}\mu_t(\dd q)$, $z\geq0$,
\begin{equation}
 \partial_t\phi_t(z)
 =2b[\phi_t(z/2)-\phi_t(z)]
       +b[1-\phi_t(z)]\partial_z\phi_t(z),
 \label{eq:critical-Laplace-PDE}
\end{equation}
with the right derivative in $z$ at zero, and with an independent $S$
from \eqref{eq:critical-product},
\begin{equation}
 h(t)=1-\E_S\phi_t(S),\qquad
 -h'(t)/b=\E_S[-\partial_z\phi_t(S)\,\phi_t(S)].
 \label{eq:Laplace-mixture}
\end{equation}
All these formulas hold in the finite-count, finite-mass solution
class without higher size moments.
\end{theorem}

Equation \eqref{eq:critical-Laplace-PDE} is the classical
Laplace-transform reduction of the additive-kernel Smoluchowski
equation \citep{Golovin1963,Bertoin2002} with the dilation term
$2b[\phi_t(z/2)-\phi_t(z)]$ added by equal splitting; it is recorded
for the mixture identities, not as a new equation.

\begin{proof}
The conditional no-future-event probability is $\Psi(q)$ by
\cref{thm:daughter-envelopes}, so averaging proves
\eqref{eq:exact-last-observable}. Apply \eqref{eq:critical-Q-generator}
to the bounded test $\Psi$; the recurrence in
\eqref{eq:product-identities} cancels the fragmentation term against
the loss part of the coagulation term:
\[
 \mathcal G_t\Psi(q)
 =2b[\Psi(q/2)-\Psi(q)]
       +bq\int_E[\Psi(q+v)-\Psi(q)]\mu_t(\dd v)
 =bq\int_E\Psi(q+v)\mu_t(\dd v).
\]
By $\Psi(q)\leq(1+q)^{-1}$, $q\Psi(q+v)\leq q\Psi(q)\leq1$, and the
coagulation term before cancellation is bounded by $b$. The
bounded-test forward equation \eqref{eq:number-forward} therefore
yields, for $s\leq t$,
\begin{equation}
 h(t)-h(s)=-b\int_s^t\iint_{E^2}q\Psi(q+v)
                         \mu_u(\dd q)\mu_u(\dd v)\dd u.
 \label{eq:exact-density-integrated}
\end{equation}
The solution is continuous in total variation by \cref{def:solution},
and the constant normalization at criticality makes $\mu_t$, hence its
product law, continuous in total variation; the bounded integrand in
\eqref{eq:exact-density-integrated} thus has an expectation continuous
on $[0,\infty)$, proving the derivative assertion. Since $L<\infty$
almost surely, this derivative and $1-h(0)=\E\Psi(Q_0)$ exhaust the
law of $L$.

The first size moment gives $\partial_z\phi_t(z)=-\E[Q_te^{-zQ_t}]$,
including the right derivative $-1$ at zero. At fixed $z>0$ both
$e^{-zq}$ and $qe^{-zq}$ are bounded; the coagulation contribution to
the bounded-test equation for $\phi$ is
$b\E[Q_te^{-zQ_t}](\phi_t(z)-1)=b[1-\phi_t(z)]\partial_z\phi_t(z)$,
and fragmentation contributes $2b[\phi_t(z/2)-\phi_t(z)]$. Continuity
of these bounded expectations gives the pointwise time derivative in
\eqref{eq:critical-Laplace-PDE}; at $z=0$ both sides vanish. Finally
the Laplace representation of $\Psi$, independence of $Q_t,Q'_t,S$,
and Tonelli give $\E\Psi(Q_t)=\E_S\phi_t(S)$ and
$\E[Q_t\Psi(Q_t+Q'_t)]=\E_S\{\E[Q_te^{-SQ_t}]\E[e^{-SQ'_t}]\}$, which
is \eqref{eq:Laplace-mixture}.
\end{proof}

The density counts an event followed by no later coagulation; the
expected rate of all coagulations remains $b$. The scalar mixture
\eqref{eq:Laplace-mixture} does not by itself give an evolution
equation in $h$ alone.

\begin{theorem}[Last-event density for general critical daughters]
\label{thm:general-last-density}
Allow time- and parent-dependent daughters. With the restarted
survival function $w_t$ of \cref{sec:daughter-comparison}, the law of
$L$ on $(0,\infty)$ is absolutely continuous, with density
\begin{equation}
 j(t)=b\E[Q_t w_t(Q_t+Q'_t)]
 \quad\text{for almost every }t>0,
 \label{eq:general-last-density}
\end{equation}
and atom $\E w_0(Q_0)$ at zero. For every $t\geq s\geq0$,
\begin{equation}
 h(t)\geq h(s)e^{-b(t-s)}.
 \label{eq:lower-calendar-tail}
\end{equation}
The coefficient $b$ is instantaneously sharp as a uniform constant
over initial populations, already within equal splitting.
\end{theorem}

\begin{proof}
The map $(t,q)\mapsto w_t(q)$ is jointly measurable: with $Y$ started
at $q$ at time $t$ and driven by the fragmentation Poisson measure
through the measurable quantile map,
$w_t(q)=\E\exp(-\int_0^\infty bY_s\dd s)$ is the expectation of a
composition of measurable maps of $(t,q)$ and the Poisson points,
obtained as the monotone limit of finite-horizon survival
probabilities. Enumerate the auxiliary coagulation times as $\tau_i$.

Two steps give the density. (i) By the strong Markov property at the
stopping time $\tau_i$, the Poisson increments after $\tau_i$ are
independent of $\mathcal F_{\tau_i}$, so the conditional probability
given $\mathcal F_{\tau_i}$ that no coagulation follows $\tau_i$ is
$w_{\tau_i}(Q_{\tau_i})$, evaluated at the post-event state; this
conditions at a genuine stopping time, not at $L$. (ii) For a Borel
set $A\subset(0,T]$, pathwise exactly one coagulation has no successor
when $L\in A$ and none otherwise, so
$\Prob(L\in A)=\E\sum_i\1_{\{\tau_i\in A\}}w_{\tau_i}(Q_{\tau_i})$,
an integrable sum since the expected number of coagulations on
$(0,T]$ is $bT$. The marked coagulation point process has compensator
$bQ_{t-}\mu_t(\dd v)\dd t$ on the partner mark $v$, and the integrand
$w_t(Q_{t-}+v)$ is predictable in $t$, so
\[
 \Prob(L\in A)
 =\int_A b\iint q\,w_t(q+v)\mu_t(\dd q)\mu_t(\dd v)\dd t,
\]
where the pre-event marginal is $\mu_t$ for almost every $t$ because
the prescribed jump intensities have no fixed-time atoms. This proves
\eqref{eq:general-last-density}, local absolute continuity of $h$, and
the atom at zero by the definition of $w_0$.

The critical envelopes \eqref{eq:critical-envelopes} imply
\[
 j(t)\leq b\E\frac{Q_t}{1+Q_t+Q'_t}
      \leq b\E\frac{Q_t}{1+Q_t}\leq b\E g_t(Q_t)=bh(t).
\]
Hence $h'=-j\geq-bh$ almost everywhere; multiplication by $e^{bt}$ and
absolute continuity give \eqref{eq:lower-calendar-tail}.

For sharpness take equal splitting and
$Q_0\sim\nu_{\varepsilon^2,\varepsilon}$ from
\eqref{eq:two-point-family}, with large atom
$r=r(\varepsilon^2,\varepsilon)$. Since $g'(0)=1$,
$h(0)=(1-\varepsilon^2)g(\varepsilon)+\varepsilon^2g(r)\sim\varepsilon$,
while the pair of small states gives
$\E[Q_0\Psi(Q_0+Q'_0)]\geq(1-\varepsilon^2)^2\varepsilon\Psi(2\varepsilon)
\sim\varepsilon$. The proved upper bound is $h(0)$, so
$-h'(0)/(bh(0))\to1$: any coefficient smaller than $b$ fails for some
such initial law, using the right derivative at zero.
\end{proof}

\begin{corollary}[Exponential-moment bracket at criticality]
\label{cor:critical-moment-bracket}
For every allowed critical daughter kernel,
\begin{equation}
 h(0)e^{-bt}\leq h(t)\leq
       \mathcal A_{1/2}(0)e^{-\kappa bt},\qquad
       \kappa=3-2\sqrt2,
 \label{eq:critical-tail-bracket}
\end{equation}
and
\begin{equation}
 \E e^{rL}<\infty\quad(0<r<\kappa b),\qquad
 \E e^{rL}=\infty\quad(r\geq b).
 \label{eq:critical-exp-bracket}
\end{equation}
\end{corollary}

\begin{proof}
Use \cref{thm:general-last-density,cor:last-moments-counts}, with
$\chi_{1/2}=2ba_{1/2}=\kappa b$. The lower envelope
$g_t(q)\geq q/(1+q)$ and $Q_0>0$ give $h(0)>0$. The identity
$\E e^{rL}=1+r\int_0^\infty e^{rt}h(t)\dd t$, allowing infinity, gives
divergence even at $r=b$.
\end{proof}

These are proved bounds; they do not identify the large-time exponent
or the exponential-moment abscissa within $[\kappa b,b]$. The next
statement reduces that question to the deterministic equation.

\begin{proposition}[Reduction of the last-event exponent to the overlap]
\label{prop:overlap-reduction}
For every critical daughter kernel, with $O(t)$ the deterministic
number--mass overlap of \eqref{eq:overlap-last-def},
$\tfrac12O(t)\leq h(t)\leq O(t)$ for all $t\geq0$. Consequently
$\liminf_{t\to\infty}(-\log h(t))/t=\liminf_{t\to\infty}(-\log O(t))/t$,
the same holds for the limits superior, and
$\sup\{r:\E e^{rL}<\infty\}=\sup\{r:\int_0^\infty e^{rt}O(t)\dd t<\infty\}$.
All these exponents lie in $[\kappa b,b]$ by
\eqref{eq:critical-tail-bracket} and \eqref{eq:overlap-affinity}.
\end{proposition}

\begin{proof}
The two-sided bound is \cref{cor:overlap-all-rates} at $\sigma=b$. A
constant factor changes neither exponential rates nor the finiteness
of $\int_0^\infty e^{rt}h(t)\dd t$, which determines $\E e^{rL}$ by
the identity in the preceding proof.
\end{proof}

The calendar-time exponent of the last auxiliary coagulation is
therefore a property of $n_t$ alone: the exponential decay rate of
$\int\min(1,N_0x/m)\,\eta_t(\dd x)$. The open question is purely
about the deterministic equation.

\begin{remark}[Conjecture for monodisperse equal-split data]
\label{rem:exponent-conjecture}
For monodisperse critical data with equal splitting we conjecture that
$O(t)$ and $h(t)$ decay with exponent $b/2$, the pure-coagulation
Golovin--Borel exponent of \eqref{eq:Borel-last-asymptotic}
\citep{Golovin1963,Bertoin2009}; the certified bracket
\eqref{eq:critical-tail-bracket} confines the exponent to
$[\kappa b,b]$ with $\kappa b\approx0.17b$. The nine retained particle
runs of \cref{sec:finite} point toward the conjectured value: the
fitted decay rates of the finite-system half moment on $[4,10]$ listed
in \cref{tab:finite-decay} are $0.23b$--$0.28b$ for $n=10^3$,
$0.30b$--$0.37b$ for $n=10^4$, and $0.43b$--$0.50b$ for $n=10^5$,
increasing with $n$ and well above $\kappa b$; an additional exploratory Marcus--Lushnikov run with
$n=2\times10^4$ and $5\times10^4$, not archived, gave log-slopes of
$O(t)$ on $[4,9]$ of $0.54b$ and $0.57b$. This is numerical evidence
subject to finite-size effects (the mass ceiling of \cref{sec:finite}
keeps $M_{1/2}\geq n^{-1/2}$ in the unit-mass runs) and to a short
fitting window; it is not a proof, and $\mathcal A_{1/2}$ bounds $O$
above rather than equalling it.
\end{remark}

\begin{remark}[A two-point obstruction to uniform scalar dissipation]
\label{rem:scalar-obstruction}
For equal splitting and a positive mean-one law $\mu$ put
\[
 \mathfrak h(\mu)=\int_E[1-\Psi(q)]\mu(\dd q),\qquad
 \mathfrak d(\mu)=\iint_{E^2}q\Psi(q+v)\mu(\dd q)\mu(\dd v),
\]
so that $h(0)=\mathfrak h(\mu_0)$ and $h'(0)=-b\,\mathfrak d(\mu_0)$
by \cref{thm:exact-last-density}. One might hope to close the bracket
\eqref{eq:critical-tail-bracket} by an instantaneous scalar inequality
$h'\leq-cb\,h^\beta$ holding uniformly over initial data. For
$\beta<1$ this is impossible for a trivial reason: it forces $h$ to
reach zero in finite time, contradicting $h(t)\geq h(0)e^{-bt}>0$. For
$\beta\geq1$, with $\beta=1$ the relevant case, a two-point family
excludes it: for $\varepsilon\downarrow0$ let
\begin{equation}
 q^-_\varepsilon=\exp[-(\log(1/\varepsilon))^3],\qquad
 q^+_\varepsilon=\frac{1-(1-\varepsilon)q^-_\varepsilon}\varepsilon,
 \qquad
 \mu_\varepsilon=\nu_{\varepsilon,q^-_\varepsilon}
 =(1-\varepsilon)\delta_{q^-_\varepsilon}
   +\varepsilon\,\delta_{q^+_\varepsilon}.
 \label{eq:two-point-counterexample}
\end{equation}
Since $q^+_\varepsilon\sim\varepsilon^{-1}$ and
$g(q^+_\varepsilon)\to1$,
$\varepsilon g(q^+_\varepsilon)\leq\mathfrak h(\mu_\varepsilon)
\leq\varepsilon+q^-_\varepsilon$, so
$\mathfrak h(\mu_\varepsilon)\sim\varepsilon$. Separating the first
coordinate into its two values, monotonicity of $\Psi$ and
$\varepsilon q^+_\varepsilon\leq1$ give
$\mathfrak d(\mu_\varepsilon)\leq q^-_\varepsilon
+\Psi(q^+_\varepsilon)$. The first $k$ product factors give, for every
integer $k\geq1$,
\begin{equation}
 \Psi(R)\leq 2^{k(k+1)/2}R^{-k},
 \label{eq:product-superpower}
\end{equation}
so for fixed $\beta$ and $k>\beta$ both
$q^-_\varepsilon/\varepsilon^\beta$ and
$\Psi(q^+_\varepsilon)/\varepsilon^\beta$ tend to zero:
\begin{equation}
 \inf_{\substack{\mu\text{ supported on two positive points}\\
                  \int q\mu(\dd q)=1}}
      \frac{\mathfrak d(\mu)}{\mathfrak h(\mu)^\beta}=0
 \qquad(\beta>0).
 \label{eq:scalar-obstruction}
\end{equation}
Realizing $\mu_\varepsilon$ as initial data by the scaling after
\cref{cor:overlap-all-rates}, and using continuity of $h'$ to extend a
strict violation at zero to a positive time interval, refutes every
uniform inequality $h'\leq-cb\,h^\beta$. By the large-$q$ form
\eqref{eq:product-periodic} the computation excludes more than powers:
along this family
$\mathfrak d\leq\exp[-(1+o(1))(\log(1/\mathfrak h))^2/(2\log2)]$.

The mechanism is a deferred huge atom: almost all number sits at a
size whose own coagulation hazard is negligible, while the rare huge
particle will certainly coagulate, but not yet. The family refutes
only uniform-in-initial-data instantaneous bounds of the form
$h'\leq-cb\,h^\beta$; it says nothing about data-dependent or
asymptotic scalar bounds, and it contradicts neither the upper tail
certificate, whose prefactor carries the separate information
$\mathcal A_{1/2}(0)$, nor the conjecture above. Every
$\mu_\varepsilon$ has finite moments of every order. The small atom
$q^-_\varepsilon$ decays faster than every power of $\varepsilon$ yet
stays representable in double precision (about $10^{-232}$ at
$\varepsilon=3\times10^{-4}$), which keeps the illustrative
floating-point evaluations in
\texttt{supplement/last\_coagulation\_exact.py} inside the
positive-size model; $e^{-1/\varepsilon}$ would underflow to zero
there.
\end{remark}
